\documentclass[8pt,a4paper]{article}
\usepackage[utf8]{inputenc}
\usepackage[spanish]{babel}
\usepackage{graphicx}
\usepackage{amsmath}
\usepackage{listings}
\usepackage{xcolor}
\usepackage{geometry}
\usepackage{hyperref}
\usepackage{multicol}

\geometry{margin= 2.5cm}

% Configuración para código Python
\lstset{
    language=Python,
    basicstyle=\ttfamily\tiny,
    keywordstyle=\color{blue},
    commentstyle=\color{green!60!black},
    stringstyle=\color{red},
    numbers=left,
    numberstyle=\tiny\color{gray},
    breaklines=true,
    frame=single,
    backgroundcolor=\color{gray!10}
}

\usepackage{fontspec}
\setmainfont{EB Garamond}[
    UprightFont = * Regular,
    ItalicFont = * Italic,
    BoldFont = * SemiBold,
    BoldItalicFont = * SemiBold Italic,
]


% --- Espaciado ---
\usepackage{setspace}
\setstretch{1.2}
\setlength{\parskip}{0.2em}

\begin{document}

% Portada personalizada
\begin{titlepage}
    \centering
    \vspace*{1cm}
    
    {\LARGE\textbf{Universidad de las Américas Puebla}}\\[0.5cm]
    {\large UDLAP}\\[2cm]
    
    {\Huge\textbf{Clasificación de Imágenes según Corrientes Artísticas}}\\[0.8cm]
    
    {\large\textbf{Autores:}}\\[0.3cm]
    {\large
        Erick Walter Trejo Durán\\
        Heriberto Espino Montelongo\\
        Nuria Arroyo Bustamante
    }\\[2cm]
    
    {\large\textbf{Materia:}}\\[0.2cm]
    {\large Redes Neuronales}\\[1cm]
    
    {\large\textbf{Profesor:}}\\[0.2cm]
    {\large Héctor Saib Maravillo Gómez}\\[2cm]
    
    \vfill
    
    {\large \today}
    
\end{titlepage}

\newpage

\section{Objetivo}

El objetivo es crear un sistema que pueda clasificar automáticamente pinturas en sus corrientes artísticas usando redes neuronales. Queremos que la computadora aprenda a reconocer si una obra es Impresionista, Cubista, Expresionista Abstracta, Barroca o Surrealista. Trabajamos con un dataset de 200 imágenes distribuidas en estos 5 géneros: Impresionismo (75), Surrealismo (50), Barroco (30), Cubismo (30) y Expresionismo abstracto (15).

\section{Metodología}

\subsection{Preprocesamiento y Extracción de Características}

Primero preparamos las imágenes normalizando los valores de píxeles de [0, 255] a [0, 1] y redimensionándolas a un tamaño estándar. Luego extraemos las diferentes características por imagen que capturan diferentes aspectos visuales:

\textbf{Color RGB:} Media, desviación estándar, mínimos y máximos de cada canal. El Barroco tiende a paletas oscuras, mientras que el Impresionismo es más claro.

\textbf{Saturación y brillo (HSV):} Obras con saturación alta (Expresionismo abstracto) son muy intensas, mientras que paletas apagadas tienen valores bajos.

\textbf{Contraste (RMS):} El Barroco con su tenebrismo tiene contraste alto, mientras que el Impresionismo es más suave.

\textbf{Textura:} Filtros Sobel (bordes), Laplaciano (detalles finos) y Gabor en 4 orientaciones capturan el estilo de pincelada. El Cubismo tiene bordes definidos, el Impresionismo es difuso.

\textbf{Composición luminosa:} Dividimos la imagen en rejilla 3×3 y analizamos dónde está la luz. El Barroco tiene foco central luminoso con fondo oscuro; el Impresionismo tiene luz homogénea.

\textbf{Entropía:} Mide el caos visual. Pollock tiene entropía alta, estilos controlados tienen valores bajos.

\textbf{LBP:} Captura textura microscópica de la pincelada, distinguiendo superficies lisas de empastes marcados.

Al final, cada imagen queda como vector $x\in\mathbb{R}^d$ donde $d\approx70$.

\subsection{Manejo del Desbalanceo}

Debido a que el dataset está desbalanceado. se empleó un método de sobremuestreo sintético, en el caso de la RBF se utiliza \emph{SMOTE} para generar ejemplos de clases minoritarias. Para la MLP se calcula peso balanceado por clase: $w_c=\frac{N}{P\,n_c}$ donde $N$ es total de muestras, $P$ número de clases y $n_c$ muestras de clase $c$. Esto solo se hace en entrenamiento; los datos de prueba se mantienen originales.

\subsection{Arquitectura de las Redes Neuronales}

\textbf{MLP (Multi-Layer Perceptron):} Red con capas totalmente conectadas. Arquitectura estrecha para controlar sobreajuste:
\[
x \xrightarrow{\text{Linear}(d\to h)} \text{BN} \xrightarrow{\text{ReLU}} \text{Dropout}(p)
\xrightarrow{\text{Linear}(h\to b)} \tanh \xrightarrow{\text{Dropout}(p)}
\xrightarrow{\text{Linear}(b\to P)} z
\]

Donde $h$ es el ancho de primera capa, $b$ es el cuello de botella (dimensión latente compacta que fuerza a aprender factores globales), $p$ es probabilidad de dropout, y $P=5$ clases. El cuello de botella reduce capacidad ($b\ll h$) para evitar sobreajuste. Usamos BatchNorm para estabilidad y dropout alto para regularización.

\textbf{Función de pérdida:} Entropía cruzada ponderada $\mathcal{L}_{\text{WCE}}(z,y)= - w_{y}\,\log(\text{softmax}(z)_y)$. También probamos label smoothing y focal loss.

\textbf{Optimización:} AdamW con learning rate y weight decay como hiperparámetros. Early stopping monitoreando macro-F1 en validación: $\hat{\theta}=\arg\max_{\theta} \text{F1}_{\text{macro,val}}(\theta)$.

\textbf{Validación cruzada:} StratifiedKFold con $K=5$ folds para preservar proporciones de clase. Para cada configuración de hiperparámetros promediamos: $S(\text{hp})=\frac{1}{K}\sum_{k=1}^K \text{F1}_{\text{macro},k}$.

\textbf{Búsqueda de hiperparámetros:} RandomizedSearchCV explorando: \texttt{hidden\_first}$\in\{16,32,48,64\}$, \texttt{bottleneck}$\in\{8,12,16,24\}$, $p\in\{0.2,\dots,0.6\}$, lr$\sim\log U(10^{-4},3\cdot10^{-3})$, weight\_decay$\sim\log U(10^{-6},3\cdot10^{-4})$.

El mejor MLP alcanzó accuracy de 0.735 ± 0.03 en validación cruzada en la clasificación por artistas.

\textbf{SVM con kernel RBF:} Clasificador de vectores de soporte con kernel de base radial, muy efectivo en espacios de alta dimensión. Optimizamos hiperparámetros C y $\gamma$ con GridSearchCV (más eficiente que RandomSearch al tener solo 2 parámetros). Este modelo fue el más estable, con accuracy de 0.785 ± 0.02 en validación cruzada.

\textbf{Bagging:} Entrenamos varios modelos independientes cambiando la semilla. En inferencia promediamos logits: $\bar{z}(x)=\frac{1}{M}\sum_{m=1}^M z_m(x)$, $\hat{y}(x)=\arg\max_c \bar{z}_c(x)$. Esto reduce varianza y mejora generalización.

\section{Resultados}

Comparamos tres modelos finales: (1) SVM RBF balanceado, (2) Bagging MLP optimizado, y (3) Bagging SVM.

El \textbf{SVM RBF balanceado} mostró el mejor equilibrio entre desempeño y estabilidad. Obteniendo un accuracy de 0.82 y un F1-macro de 0.80 en el conjunto de prueba, indicando alta consistencia del modelo. Los géneros más fáciles de clasificar fueron Barroco, Impresionismo y Expresionismo abstracto (F1 $\geq$ 0.90), porque tienen características visuales muy distintivas. El Cubismo fue el más difícil, porque comparte elementos geométricos con otros estilos.

La validación cruzada estratificada confirmó consistencia (std = 0.02), indicando que el modelo no depende de un split afortunado. Esto es crucial con solo 200 imágenes.

El \textbf{Bagging SVM} alcanzó la mejor métrica puntual en prueba (accuracy $\approx$ 0.85, F1-macro $\approx$ 0.86), pero la validación cruzada reveló mayor varianza (std $\approx$ 0.064) y menor desempeño promedio, sugiriendo sensibilidad al muestreo y posible sobreajuste al split de prueba.

Los \textbf{Bagging MLP Optimizado} obtuvieron métricas con accuracy $\approx$ 0.70 y F1-macro $\approx$ 0.70 para la clasificación por pintor, pero mostraban mayor inestabilidad (std $\approx$ 0.03) y tendencia a sobreajustar, probablemente por la complejidad del modelo y el tamaño limitado del dataset. Se aplicó bagging y, luego de traducir a corriente artística, se logró accuracy $\approx$ 0.95 y F1-macro $\approx$ 0.956.

\section{Evaluación}

Usamos Accuracy (porcentaje de predicciones correctas) y F1-Score (balancea precisión y recall, útil con clases desbalanceadas). Las matrices de confusión mostraron que el modelo a veces confunde Cubismo con otros estilos por similitudes visuales de la obra.

\section{Conclusiones}

El proyecto demuestra que es posible clasificar corrientes artísticas con precisión razonable. El \textbf{Bagging MLP Optimizado} es el modelo final seleccionado por ofrecer el mejor equilibrio entre rendimiento (95\% accuracy en 5 clases, azar daría 20\%), estabilidad y capacidad de generalización.

Los géneros con características distintivas (Barroco con claroscuro dramático, Impresionismo con colores vivos) son más fáciles de clasificar. El Cubismo presenta desafíos por similitudes con otros estilos.

Las limitaciones principales son el tamaño del dataset (200 imágenes) y que usamos features hechas a mano. El modelo parece haber alcanzado un "techo" natural con estos datos.

Mejoras futuras: (1) más imágenes, (2) CNNs pre-entrenadas que extraen features automáticamente.

\end{document}